% Regenerated by regen_tables.py. Ledger L-2026-09-01-08, -22 (T0-5).
\begin{table}[h]
\centering
\caption{Four-point hole reorganization energy for every molecule in the charged-state subset. Molecule 19531 is a pyrrolidinium ester chloride ion pair, rejected by the validity filter as a multi-fragment species: removing an electron redistributes charge between its fragments rather than ionising a neutral molecule, so the four-point expression does not describe a reorganization and its value is not comparable with the others. Its $\lhole$ is nevertheless an order of magnitude more negative than either azido-triazine, which shows that a negative $\lhole$ in this subset is not specific to the azide group. Molecules 14380 and 18985 are the two remaining members of the same convenience batch, and their gaps of \qtylist{5.23;5.37}{\electronvolt} place both far outside the photovoltaic window, so they fix the sign of $\lhole$ for ordinary closed-shell molecules and support nothing beyond that.}
\label{tab:reorganization}
\begin{tabular}{lccccl}
\toprule
ID & \makecell{$\lhole$ B3LYP\\(\unit{\electronvolt})} & \makecell{CAM-B3LYP\\(\unit{\electronvolt})} & \makecell{$\omega$B97X-D3\\(\unit{\electronvolt})} & SCF status & note \\
\midrule
19531 & $-3.444$ & $-4.013$ & $-4.209$ & \textbf{anion\_neutral\_failed} & \textbf{chloride ion pair}; formula not applicable \\
20778 & $-0.508$ & $-0.812$ & $-0.903$ & converged & azide \\
17851 & $-0.252$ & $-0.472$ & $-0.530$ & converged & azide \\
18985 & $+0.406$ & $+0.269$ & $+0.240$ & converged & wide gap (\qty{5.37}{\electronvolt}); pyridyl imine \\
977 & $+0.548$ & $+0.558$ & $+0.561$ & converged & \ch{O2}, not a molecule of interest \\
14380 & $+0.949$ & $+0.817$ & $+0.715$ & converged & wide gap (\qty{5.23}{\electronvolt}); sulfonamide \\
\bottomrule
\end{tabular}
\end{table}
